\section{嘉宾介绍与访谈背景}

本期 WhynotTV Podcast（第三期）的嘉宾是陈天奇（Tianqi Chen），CMU 机器学习系助理教授，OctoML 联合创始人（2024 年底被 NVIDIA 收购）。过去近 20 年，陈天奇围绕\textbf{``让机器学习更轻、更快、更易部署''}这条主线，先后主导了 XGBoost、MXNet、TVM、MLC LLM 等广受欢迎的开源项目。其中 XGBoost 论文引用数接近 7 万次，至今仍是 tabular data 领域最常用的工具。

2019 年他在知乎发表《机器学习科研这十年》，收获上万点赞。六年后的 2025 年，他在机器学习这条路上已经走了近 20 年。主持人（WhynotTV）指出：``对我而言，陈天奇不仅仅是 AI 领域一个响亮的名字，他还影响了我对机器学习和科研最早的认知。''

本次访谈长达 160 分钟，从他的童年经历出发，经过本科 ACM 班、研究生 APEX 实验室、UW 博士期间的四大项目（XGBoost、MXNet、TVM、Google Brain 实习的 Net2Net），到 CMU 教职和 OctoML 创业，最后上升到长期主义、初心和勇气等人生哲学。

\begin{importantbox}{陈天奇的技术主线}
从 XGBoost（2014）到 MXNet（2015）到 TVM（2017）再到 MLC LLM（2023），陈天奇的每一个项目都围绕同一个核心问题：\textbf{如何用更少的工程资源、更自动化的方式，让机器学习系统在各种硬件上高效运行}。这条主线贯穿了大数据时代、深度学习时代和大模型时代。他不是一个绑定在某种方法上的研究者，而是一个追踪问题的人------方法可以变（从树模型到深度学习框架到编译器），但问题始终是``让 ML 系统更好''。
\end{importantbox}

